%% =====================================================================
%% CANONICAL TEMPLATE — arXiv / preprint (single-column, article class)
%% Copy this file, keep the structure, swap the content.
%% Build:  writter-ai/skills/academic-paper/scripts/build.sh example-arxiv.tex
%% arXiv accepts standard `article`. To match a venue, drop in its style file
%% (e.g. neurips_2024.sty, icml2024.sty) and replace the \usepackage line noted below.
%% =====================================================================
\documentclass[11pt]{article}

\usepackage[letterpaper,margin=1in]{geometry}   % venue style files usually set this
\usepackage[utf8]{inputenc}
\usepackage[T1]{fontenc}
\usepackage{lmodern}
\usepackage{amsmath,amssymb,amsfonts,amsthm}
\usepackage{graphicx}
\usepackage{booktabs}
\usepackage{authblk}                 % clean author/affiliation blocks
\usepackage[numbers,sort&compress]{natbib}   % [1,2,4] numeric; or drop `numbers` for (Author, Year)
\usepackage{xcolor}
\usepackage[hidelinks,breaklinks=true]{hyperref}
% ---- To use a venue style instead, comment geometry+authblk above and add e.g.:
% \usepackage[preprint]{neurips_2024}

\theoremstyle{plain}
\newtheorem{theorem}{Theorem}
\newtheorem{lemma}{Lemma}
\theoremstyle{definition}
\newtheorem{definition}{Definition}

\title{\bfseries A Concise, Specific Paper Title}

\author[1]{First A. Author}
\author[1,2]{Second B. Author}
\author[2]{Third C. Author}
\affil[1]{Department, Organization, \texttt{\{first,second\}@org.edu}}
\affil[2]{Second Organization, \texttt{third@org2.edu}}
\date{}   % arXiv stamps its own date; keep empty

\begin{document}
\maketitle

\begin{abstract}
Problem in one sentence, gap in the second, approach in the third, headline
result in the fourth. 150--250 words, no citations, every claim proven in the
body. Replace this paragraph with yours.
\end{abstract}

\section{Introduction}
Problem, why it matters, the specific gap, and your contribution. End with an
explicit contribution list:
\begin{itemize}
\item A clearly stated method (Section~\ref{sec:method}).
\item An empirical result improving X by $Y\%$ over \citet{ref_baseline}.
\item Analysis/ablation isolating the gain (Section~\ref{sec:results}).
\end{itemize}

\section{Related Work}
Group prior art into threads; state what each does and its limitation relative to
your setting \citep{ref_survey,ref_method}.

\section{Method}\label{sec:method}
Define notation once; make the approach reproducible. A numbered equation:
\begin{equation}\label{eq:objective}
\mathcal{L}(\theta) = \frac{1}{N}\sum_{i=1}^{N}\ell\!\big(f_\theta(x_i), y_i\big) + \lambda\lVert\theta\rVert_2^2 .
\end{equation}
State results as formal claims where appropriate:
\begin{theorem}\label{thm:main}
Under assumptions A1--A2, the estimator $\hat\theta$ satisfies $\lVert\hat\theta-\theta^\star\rVert = \mathcal{O}(1/\sqrt{N})$.
\end{theorem}
\begin{proof}[Proof sketch] Give the idea; defer full proof to an appendix. \end{proof}

\section{Experiments}\label{sec:exp}
Datasets, baselines, metrics, protocol (splits, seeds, hardware). Every reported
number must be reproducible from this section.

\section{Results}\label{sec:results}
Lead with the table, then analysis.
\begin{table}[t]
\centering
\caption{Main results. Best per column in \textbf{bold}; mean$\pm$std over seeds.}
\label{tab:main}
\begin{tabular}{lcc}
\toprule
Method & Metric A ($\uparrow$) & Metric B ($\downarrow$) \\
\midrule
Baseline \citep{ref_baseline} & 71.2 & 0.184 \\
Ours                          & \textbf{78.6} & \textbf{0.142} \\
\bottomrule
\end{tabular}
\end{table}

%\begin{figure}[t]
%  \centering
%  \includegraphics[width=0.8\linewidth]{figure.pdf}
%  \caption{Caption states the takeaway, not just the contents.}
%  \label{fig:overview}
%\end{figure}

\section{Conclusion}
Restate the contribution and the key number; one limitation, one next step.

\section*{Acknowledgments}
Optional: funding, compute, non-author contributors.

\bibliographystyle{plainnat}   % or abbrvnat / unsrtnat; venue styles override this
\bibliography{refs}

\appendix
\section{Full Proof of Theorem~\ref{thm:main}}
Defer long derivations here.

\end{document}
